\documentclass[10pt,a4paper]{amsart}
\usepackage[margin=19mm]{geometry}
\usepackage{amsmath,amssymb,mathtools}
\usepackage[scr=rsfs,cal=euler]{mathalfa}
\usepackage{xfrac}
\usepackage[unicode,hidelinks,bookmarks=false]{hyperref}
\newcommand{\ie}{i.e.\ }
\newcommand{\cf}{cf.\ }
\newcommand{\resp}{respectively\ }
\newcommand{\Subsection}{\subsection}
\providecommand{\nobreakdash}{\nobreak\hbox{-}}
\setlength{\emergencystretch}{2em}
\multlinegap0pt
\usepackage{amssymb}
\usepackage{booktabs}
\usepackage[shortlabels]{enumitem}
\usepackage{tikz}
\newtheorem{theorem}{Theorem}
\title{Group-Theoretical Origin of the Sectoral-Tesseral-Zonal Trichotomy in Spherical Harmonics}
\author{Mustafa Bakr, Smain Amari}
\date{Source: arXiv:2512.20119v1 (CC BY 4.0)}
\begin{document}
\maketitle

\noindent\emph{Source and licence.} Group-Theoretical Origin of the Sectoral-Tesseral-Zonal Trichotomy in Spherical Harmonics. Version arXiv:2512.20119v1 (CC BY 4.0), distributed by arXiv under the Creative Commons Attribution 4.0 International licence, http://creativecommons.org/licenses/by/4.0/, as recorded in the arXiv licence metadata for this exact version and shown on the abstract page https://arxiv.org/abs/2512.20119v1. Full licence notice: \texttt{licenses/arxiv-2512.20119-CC-BY-4.0.txt}.\par\smallskip

\noindent\textit{Editorial scope.} Counted unit: the article's theorem thm:sectoral (source0.tex lines 117--125). The article's complete original proof of it is delivered with it (source0.tex lines 127--129, one \texttt{proof} environment of 3 lines). No mathematics is written, repaired or re-derived: the statement and the proof are the article's own lines, in the article's own notation, and no retained line is substituted or re-flowed.  No further source-local context is needed: every cross-reference of every retained line resolves inside the two delivered spans.  Nothing was omitted from any retained span, so the package carries no omission record.  No retained line contains a citation, so the wrapper carries no bibliography and no bibliography entry.  No retained line contains a figure environment, an image-inclusion command or drawing code, so no image is delivered and the package declares no asset dependency.  Editorial formatting only, in addition: an AMS \texttt{amsart} preamble that restates the article's own declarations and macros used by the retained lines, namely the environments \texttt{theorem} on separate counters, together with the standard packages those lines need.  The retained environments print in this document as Theorem 1 in order of appearance, whereas the article prints the same items with the numbering of its own full document.  The complete native source of the article as distributed in the arXiv e-print for this exact version is bundled unchanged as \texttt{sources/191538/source0.tex}.
\begin{theorem}[Continuous Sectoral Family]
\label{thm:sectoral}
The function
\begin{equation}
\psi_m(\theta, \phi) = (\sin\theta)^m e^{im\phi}
\label{eq:full_sectoral}
\end{equation}
satisfies $L_+ \psi_m = 0$ and $L_z \psi_m = m\psi_m$ for any real $m > 0$, and is regular at both poles of the sphere. The Casimir eigenvalue is $\mathrm{L}^2 \psi_m = m(m+1)\psi_m$.
\end{theorem}

\begin{proof}
The properties $L_+ \psi_m = 0$ and $L_z \psi_m = m\psi_m$ follow from the construction above. For the Casimir eigenvalue, we use the identity $\mathrm{L}^2 = L_- L_+ + L_z^2 + L_z$, which follows from the commutation relations. Applying this to a state annihilated by $L_+$ gives $\mathrm{L}^2 \psi_m = (L_z^2 + L_z)\psi_m = (m^2 + m)\psi_m = m(m+1)\psi_m$. Regularity at the poles follows from the behaviour $(\sin\theta)^m \to 0$ as $\theta \to 0$ or $\theta \to \pi$ for $m > 0$.
\end{proof}

\end{document}
